% Volume II: Causal Explosion
% Part IV. Time, Delay, and the Maturation of Evidence
% Chapter 19: Hierarchies of Salience
% Spine: proof-spines/ch019-spine.md

\chapter{Hierarchies of Salience}
\label{ch:ch019}

A crucial detail often must survive a hierarchy before it can matter in judgment:
\[
E_0\to E_1\to \cdots\to E_L.
\]
\Definition\ This chain defines the chapter's levels of evidentiary transmission.
If at each level \(i\) it survives with probability \(p_i\), then the probability that it reaches the final decision-maker is
\[
P_{\mathrm{reach}}
=
\prod_{i=0}^{L-1}p_i.
\]
\ToyModel\ This reach probability is a stylized filter model rather than a fitted measurement law.
This is the formal core of \textbf{hierarchies of salience}: visibility is filtered step by step rather than given all at once.

The consequence is that important evidence may remain practically invisible even when it is public. Filters, triage, editorial selection, bureaucratic referral, algorithmic ranking, and expert validation all alter the \(p_i\)'s, sometimes elevating spectacle and sometimes burying what is most probative. The chapter studies how institutions can redesign those filters so that salience becomes a tool of disciplined suspicion rather than a conveyor belt for distortion.
